\chapter{Arquitectura del sistema}
\label{ch:arquitectura}

Este capítulo describe la organización estructural del simulador: los componentes que lo forman, cómo se relacionan entre sí y las decisiones de diseño que condicionan esa organización. El detalle de implementación de cada componente se desarrolla en el Capítulo~\ref{ch:implementacion}.

El sistema está compuesto por cinco bloques funcionales: una interfaz web, un backend de orquestación, tres instancias del motor de enrutado viario OSRM, un planificador de transporte público basado en OpenTripPlanner y un módulo de inferencia de elección modal. Todos ellos se despliegan como contenedores Docker orquestados mediante Docker Compose.

% -------------------------------------------------------
\section{Visión general}
\label{sec:arch_vision_general}
% -------------------------------------------------------

La Figura~\ref{fig:arch_general} muestra la arquitectura de alto nivel del sistema. El principio de diseño central es que el frontend nunca se comunica directamente con ningún servicio externo: toda petición pasa por el backend FastAPI, que actúa como única puerta de entrada y punto de orquestación. Esta decisión simplifica el control de versiones de la API, centraliza el manejo de errores y evita exponer los puertos internos de OSRM y OTP al navegador.

\begin{figure}[htbp]
  \centering
  \includegraphics[width=\textwidth]{figs/ARQUITECTURA_pytorch_final.pdf}
  \caption{Arquitectura general del simulador de movilidad urbana.}
  \label{fig:arch_general}
\end{figure}

El frontend expone la interfaz cartográfica al usuario y delega en el backend el cálculo de rutas, la consulta de horarios y la inferencia modal. El backend integra cuatro routers diferenciados: \texttt{/api/osrm} para rutas viarias por perfil de transporte, \texttt{/api/otp} para itinerarios de transporte público, \texttt{/api/gtfs} para información estática de la red de autobús, y \texttt{/api/lpmc} para la inferencia de elección modal. Cada router encapsula la lógica de comunicación con su servicio subyacente y expone al frontend una interfaz homogénea, independiente de los detalles de cada protocolo.

% -------------------------------------------------------
\section{Infraestructura y orquestación}
\label{sec:arch_infra}
% -------------------------------------------------------

La infraestructura del sistema se define íntegramente en un fichero \texttt{docker-compose.yml} en la raíz del repositorio. El despliegue completo se obtiene con un único comando (\texttt{docker compose up -{}-build}), que elimina dependencias de entorno en la máquina anfitriona y asegura reproducibilidad entre equipos de desarrollo y producción.

Varios servicios requieren artefactos pesados generados a partir de los datos de origen: los grafos viarios de OSRM, el grafo multimodal de OpenTripPlanner y los modelos de aprendizaje automático entrenados. Estos ficheros se versionan mediante \textit{Git LFS} (\textit{Large File Storage}), una extensión de Git que almacena los ficheros de gran tamaño fuera del historial principal y los sustituye en el repositorio por punteros ligeros; de este modo, un clon del repositorio dispone de todos los artefactos sin necesidad de regenerarlos. Los ficheros almacenados en LFS son el extracto OSM de Castilla-La Mancha (\textasciitilde{}97~MB), el grafo OTP (\textasciitilde{}121~MB), el feed GTFS comprimido (\textasciitilde{}14~MB) y los modelos entrenados (XGBoost \textasciitilde{}17~MB, Random Forest \textasciitilde{}597~MB, red neuronal profunda \textasciitilde{}66~KB).

Para cubrir el caso de un clon sin LFS o de una actualización de los datos de origen, la orquestación incorpora tres servicios de inicialización idempotentes que se ejecutan una sola vez antes de los servicios principales: \texttt{gtfs-init} descomprime el feed GTFS, \texttt{osrm-setup} compila los grafos viarios OSRM por perfil si no están presentes, y \texttt{otp-build} construye el grafo de OTP si no existe. El procedimiento completo, tanto automático como manual, se describe en el Anexo~\ref{anx:despliegue}.

La Tabla~\ref{tab:servicios_docker} resume los servicios definidos, sus puertos expuestos y su función en el sistema. La Figura~\ref{fig:arch_docker} ilustra la topología de servicios, sus dependencias y los volúmenes compartidos.

\begin{table}[htbp]
  \caption{Servicios Docker Compose del simulador.}
  \label{tab:servicios_docker}
  \centering
  \begin{tabular}{llp{6.5cm}}
    \hline
    \textbf{Servicio} & \textbf{Puerto} & \textbf{Función} \\
    \hline
    \texttt{frontend}   & 5173 & Interfaz web React servida por Vite. \\
    \texttt{backend}    & 8000 & API FastAPI, orquestación y lógica de negocio. \\
    \texttt{osrm-car}   & 5001 & Motor OSRM con perfil de vehículo a motor. \\
    \texttt{osrm-bike}  & 5002 & Motor OSRM con perfil de bicicleta. \\
    \texttt{osrm-foot}  & 5003 & Motor OSRM con perfil peatonal. \\
    \texttt{otp}        & 8080 & OpenTripPlanner 2.x con grafo multimodal de Toledo. \\
    \hline
    \texttt{gtfs-init}  & — & Descomprime el feed GTFS del archivo ZIP. \\
    \texttt{osrm-setup} & — & Compila los grafos viarios OSRM por perfil si no están presentes. \\
    \texttt{otp-build}  & — & Construye el grafo multimodal de OTP si no está presente. \\
    \hline
  \end{tabular}
\end{table}

\begin{figure}[htbp]
  \centering
  \missingfigure{Diagrama de servicios Docker Compose: servicios de aplicación (frontend :5173, backend :8000, osrm-car :5001, osrm-bike :5002, osrm-foot :5003, otp :8080) y servicios de inicialización (gtfs-init, osrm-setup, otp-build, sin puerto expuesto). Dependencias entre servicios (service\_completed\_successfully) y volúmenes montados en los contenedores OSRM y OTP.}
  \caption{Servicios Docker Compose del simulador: dependencias y puertos expuestos.}
  \label{fig:arch_docker}
\end{figure}

La razón por la que OSRM se instancia tres veces en lugar de usar un único servidor multiperfil es que la imagen oficial de OSRM no admite perfiles simultáneos: cada proceso solo puede servir el grafo con el que fue preprocesado. Esta restricción motivó el cambio desde la solución inicial basada en el servidor de demostración público, que únicamente ofrecía el perfil de conducción, a la infraestructura local multiperfil descrita aquí.

% -------------------------------------------------------
\section{Backend: capa de orquestación}
\label{sec:arch_backend}
% -------------------------------------------------------

El backend está implementado con FastAPI sobre Python y se organiza en cuatro routers, cada uno responsable de un dominio funcional. La separación en routers permite desarrollar, probar y documentar cada integración de forma independiente. La Tabla~\ref{tab:endpoints_api} recoge los endpoints principales expuestos por cada router.

\begin{table}[htbp]
  \caption{Endpoints principales de la API REST del backend.}
  \label{tab:endpoints_api}
  \centering
  \begin{tabular}{llp{6.5cm}}
    \hline
    \textbf{Método} & \textbf{Ruta} & \textbf{Descripción} \\
    \hline
    POST & \texttt{/api/osrm/routes}                  & Rutas viarias por perfil de transporte (OSRM). \\
    POST & \texttt{/api/otp/routes}                   & Itinerarios de transporte público con desglose por tramos (OTP). \\
    GET  & \texttt{/api/gtfs/stops}                   & Listado de paradas con filtro por caja delimitadora. \\
    GET  & \texttt{/api/gtfs/routes}                  & Listado de líneas de la red de transporte. \\
    GET  & \texttt{/api/gtfs/routes/\{id\}}           & Detalle de línea: paradas ordenadas y trazado. \\
    GET  & \texttt{/api/gtfs/routes/\{id\}/schedule}  & Horarios de una línea por fecha. \\
    POST & \texttt{/api/lpmc/predict}                 & Inferencia de elección modal con el modelo activo. \\
    POST & \texttt{/api/lpmc/compare}                 & Inferencia con los tres modelos sobre el mismo escenario. \\
    POST & \texttt{/api/lpmc/debug-features}          & Vector de características antes y después del escalado. \\
    GET  & \texttt{/api/lpmc/models}                  & Modelos disponibles y modelo predeterminado. \\
    GET  & \texttt{/health}                           & Comprobación de disponibilidad. \\
    \hline
  \end{tabular}
\end{table}

El router \texttt{/api/osrm} recibe un par origen-destino con la lista de perfiles solicitados, consulta en paralelo las instancias OSRM correspondientes mediante \texttt{asyncio.gather} y devuelve al frontend las rutas con distancia, duración y geometría decodificada para su representación cartográfica.

El router \texttt{/api/otp} consulta OpenTripPlanner para obtener itinerarios de transporte público y gestiona la paginación entre las alternativas devueltas, hasta cinco por consulta. La respuesta incluye el desglose por tramos, con tipo de modo de transporte, geometría, paradas de inicio y fin, y agencia operadora cuando el tramo es en autobús. Los itinerarios se ordenan por duración; se selecciona automáticamente el primero que incluya al menos un tramo en autobús y, si ninguno lo incluye, el de menor duración total.

El router \texttt{/api/gtfs} expone la información estática de la red de autobús a partir del feed GTFS cargado en memoria: listado de paradas con sus líneas asociadas, listado de líneas con sus sentidos, detalle de una línea con sus paradas ordenadas y trazado, y horarios por fecha. Estos datos son independientes de OpenTripPlanner, de modo que la red puede consultarse aunque el planificador no esté disponible.

El router \texttt{/api/lpmc} recibe las coordenadas de origen y destino junto con el perfil sociodemográfico del viajero, construye el vector de características consultando OSRM y OTP internamente, aplica el escalado y ejecuta la inferencia con el modelo activo. El endpoint \texttt{/compare} ejecuta los tres modelos de forma secuencial sobre el mismo escenario y devuelve sus probabilidades para cada modo de transporte; la ejecución secuencial evita comprometer el rendimiento de la CPU compartida entre contenedores.

% -------------------------------------------------------
\section{Frontend: interfaz cartográfica}
\label{sec:arch_frontend}
% -------------------------------------------------------

El frontend está desarrollado con React~18, Vite y TypeScript. React~18 es la biblioteca de construcción de interfaces de usuario; Vite actúa como compilador y servidor de desarrollo con recarga automática del módulo; TypeScript añade comprobación estática de tipos sobre JavaScript. La representación cartográfica se realiza con React-Leaflet, que integra la biblioteca de mapas Leaflet en el modelo de componentes de React. Los iconos de la interfaz proceden de Lucide React, una biblioteca de iconos SVG de código abierto. Las peticiones asíncronas al backend se gestionan con TanStack Query, que aporta caché de respuestas, gestión de estados de carga y reintento automático.

La aplicación se estructura en dos componentes principales. \texttt{App} concentra el estado global y la lógica de negocio del cliente: gestiona los puntos de origen y destino, los modos de transporte seleccionados, los resultados de rutas, la línea activa en la red GTFS y el perfil del viajero para la inferencia. \texttt{MapView} recibe el estado relevante mediante \textit{props} —datos que el componente padre pasa al componente hijo de forma unidireccional— y se encarga de renderizar los marcadores, los trazados de ruta y las paradas de autobús sobre el mapa.

La interfaz cartográfica está complementada por una barra lateral de navegación (\textit{rail}) de ancho fijo que da acceso a los paneles funcionales del simulador: Inicio, Rutas, Red GTFS, Predicción IA, Capas y Ajustes. La capa de fondo es seleccionable entre seis opciones; el detalle de cada una se describe en la Sección~\ref{subsec:impl_frontend_capas}.

% -------------------------------------------------------
\section{Pipeline de inferencia modal}
\label{sec:arch_pipeline_inferencia}
% -------------------------------------------------------

La inferencia de elección modal requiere coordinar cuatro peticiones asíncronas para construir el vector de entrada al modelo. La Figura~\ref{fig:pipeline_inferencia} ilustra el flujo completo.

\begin{figure}[htbp]
  \centering
  \includegraphics[width=\textwidth]{figs/PIPELINE_Inferencia.pdf}
  \caption{Pipeline de inferencia de elección modal.}
  \label{fig:pipeline_inferencia}
\end{figure}

Al recibir una petición \texttt{POST /api/lpmc/predict}, el backend lanza en paralelo cuatro consultas asíncronas mediante \texttt{asyncio.gather}: una a cada una de las tres instancias OSRM (conducción, bicicleta y a pie) y una a OpenTripPlanner. Con los resultados se construye el vector de características de ruta y se añaden las variables sociodemográficas recibidas en la petición, componiendo el vector completo de entrada al modelo.

Las variables de entrada proceden de tres fuentes: los tiempos y distancias por perfil de transporte calculados por OSRM, el desglose temporal del itinerario de transporte público proporcionado por OTP, y los datos sociodemográficos y contextuales que el usuario introduce en el panel Predicción IA. El Capítulo~\ref{ch:implementacion} detalla la construcción del vector de características y las transformaciones aplicadas.

Las variables continuas se normalizan con un escalado estándar aplicado a las columnas de ruta y contextuales, preservando las variables binarias y las columnas codificadas sin transformar. El modelo devuelve un vector de cuatro probabilidades, una por modo de transporte: a pie (\textit{walk}), bicicleta (\textit{cycle}), transporte público (\textit{pt}) y conducción (\textit{drive}). El modo predicho es el de mayor probabilidad.

El modelo activo se selecciona mediante el parámetro opcional \texttt{model\_variant} de la petición; en su ausencia, se recurre a la variable de entorno \texttt{LPMC\_MODEL\_VARIANT} (\texttt{xgb} por defecto). Este mecanismo permite cambiar el modelo de inferencia y facilita la incorporación de modelos entrenados externamente, sin modificar el código del backend.
